
\documentclass{article}
\usepackage{colortbl}
\usepackage{makecell}
\usepackage{multirow}
\usepackage{supertabular}

\begin{document}

\newcounter{utterance}

\centering \large Interaction Transcript for game `cladder', experiment `full\_v1.5\_default', episode 5155 with llama{-}3.1{-}8b{-}instruct.
\vspace{24pt}

{ \footnotesize  \setcounter{utterance}{1}
\setlength{\tabcolsep}{0pt}
\begin{supertabular}{c@{$\;$}|p{.15\linewidth}@{}p{.15\linewidth}p{.15\linewidth}p{.15\linewidth}p{.15\linewidth}p{.15\linewidth}}
    \# & \multicolumn{2}{c}{Player} && \multicolumn{2}{c}{Game Master} \\
    \hline

    \theutterance \stepcounter{utterance}  
    & & & \multicolumn{4}{p{0.6\linewidth}}{
        \cellcolor[rgb]{0.9,0.9,0.9}{
            \makecell[{{p{\linewidth}}}]{
                \texttt{\tiny{[P1$\langle$GM]}}
                \texttt{You are an expert in causal inference. The following question is not a typical commonsense query, but rather a meticulously designed question created by a professor specializing in causal inference, intended to assess the students' mastery of the course content.} \\
\texttt{Imagine a self{-}contained, hypothetical world with only the following conditions, and without any unmentioned factors or causal relationships: Talent has a direct effect on elite institution admission status. Brown eyes has a direct effect on elite institution admission status. Method 1: We look at how talent correlates with brown eyes case by case according to elite institution admission status. Method 2: We look directly at how talent correlates with brown eyes in general. To understand how talent affects brown eyes, is it more correct to use the Method 1 than Method 2?} \\
\texttt{Start your answer with 'Yes' or 'No', followed by additional reasoning or evidence to support your explanation.} \\
            }
        }
    }
     \\ \\

    \theutterance \stepcounter{utterance}  
    & \multicolumn{4}{p{0.6\linewidth}}{
        \cellcolor[rgb]{0.561,0.737,0.561}{
            \makecell[{{p{\linewidth}}}]{
                \texttt{\tiny{[P1$\rangle$GM]}}
                \texttt{No, it is not more correct to use Method 1 than Method 2 to understand how talent affects brown eyes.} \\
\\ 
\texttt{In this hypothetical world, we are trying to estimate the causal effect of talent on brown eyes. However, the problem is that talent and brown eyes both have direct effects on elite institution admission status, which means that there is a common cause (elite institution admission status) that affects both variables. This is known as a "confounding variable" or "common cause" problem.} \\
\\ 
\texttt{Method 1, which looks at the correlation between talent and brown eyes case by case according to elite institution admission status, is essentially a form of "matching" or "stratification" approach. By conditioning on elite institution admission status, we are trying to control for the confounding variable and estimate the causal effect of talent on brown eyes within each stratum.} \\
\\ 
\texttt{However, this approach is not without its limitations. By conditioning on elite institution admission status, we are essentially assuming that the relationship between talent and brown eyes is the same within each stratum. But what if the relationship between talent and brown eyes is different for individuals who are admitted to elite institutions versus those who are not? This could lead to biased estimates of the causal effect.} \\
\\ 
\texttt{Method 2, which looks directly at the correlation between talent and brown eyes in general, is a more straightforward approach. However, it is also subject to bias due to the common cause problem.} \\
\\ 
\texttt{To address this problem, we would need to use a more advanced method, such as instrumental variables (IV) or regression adjustment, which can help to control for the confounding variable and estimate the causal effect of talent on brown eyes.} \\
\\ 
\texttt{In this case, neither Method 1 nor Method 2 is sufficient to estimate the causal effect of talent on brown eyes. A more advanced method is needed to address the common cause problem.} \\
            }
        }
    }
    & & \\ \\

    \theutterance \stepcounter{utterance}  
    & & & \multicolumn{2}{p{0.3\linewidth}}{
        \cellcolor[rgb]{0.9,0.9,0.9}{
            \makecell[{{p{\linewidth}}}]{
                \texttt{\tiny{[GM$|$GM]}}
                \texttt{No} \\
            }
        }
    }
    & & \\ \\

    \theutterance \stepcounter{utterance}  
    & & & \multicolumn{2}{p{0.3\linewidth}}{
        \cellcolor[rgb]{0.9,0.9,0.9}{
            \makecell[{{p{\linewidth}}}]{
                \texttt{\tiny{[GM$|$GM]}}
                \texttt{no} \\
            }
        }
    }
    & & \\ \\

    \theutterance \stepcounter{utterance}  
    & & & \multicolumn{2}{p{0.3\linewidth}}{
        \cellcolor[rgb]{0.9,0.9,0.9}{
            \makecell[{{p{\linewidth}}}]{
                \texttt{\tiny{[GM$|$GM]}}
                \texttt{game\_result = WIN} \\
            }
        }
    }
    & & \\ \\

\end{supertabular}
}

\end{document}
